\begin{examplebox}
\textbf{\textcolor{CExample}{例 1｜只知横坐标，三行求切线距离（原题纠错精简）}}

\textbf{题目：}椭圆 $\dfrac{x^2}{16}+\dfrac{y^2}{9}=1$，$P\left(2,\dfrac{3\sqrt3}{2}\right)$，过 $P$ 作切线。求原点到切线距离 $d$，并验证定理。

\textbf{解：} $a=4,b=3,c=\sqrt7$，且
\[
\frac{2^2}{16}+\frac{(3\sqrt3/2)^2}{9}=\frac14+\frac34=1,
\]
故 $P$ 确在椭圆上。焦半径直接用一次式（无需开两次平方）：
\[
r_1=4+\frac{\sqrt7}{2},\quad
r_2=4-\frac{\sqrt7}{2},\quad
r_1r_2=16-\frac74=\frac{57}{4}.
\]
由 C020，
\[
\boxed{d=\frac{ab}{\sqrt{r_1r_2}}
=\frac{12}{\sqrt{57}/2}
=\frac{24}{\sqrt{57}}}.
\]
\textbf{独立核对：}切线为 $\dfrac{x}{8}+\dfrac{\sqrt3}{6}y=1$，
\[
d=\frac1{\sqrt{1/64+1/12}}
=\frac{8\sqrt3}{\sqrt{19}}=\frac{24}{\sqrt{57}}.
\]
\textbf{方法选择：}只给 $x_0$ 时，用 $r_1r_2=a^2-c^2x_0^2/a^2$ 比开方求两段距离快。

\medskip
\textbf{\textcolor{CExample}{例 2｜反推焦半径并核查是否真有切点}}

\textbf{题目：}椭圆 $\dfrac{x^2}{25}+\dfrac{y^2}{16}=1$ 的一条切线到原点距离为 $d=2\sqrt5$。
求 $r_1r_2$，并判断可能对应的切点坐标。

\textbf{解：} $a=5,b=4,c=3$。由于 $4\le2\sqrt5\le5$，条件可实现。
\[
r_1r_2=\frac{a^2b^2}{d^2}
=\frac{25\cdot16}{20}=\boxed{20}.
\]
再用 $r_1r_2=25-\dfrac{9x_0^2}{25}$：
\[
20=25-\frac{9x_0^2}{25}
\quad\Longrightarrow\quad x_0^2=\frac{125}{9},
\quad x_0=\pm\frac{5\sqrt5}{3}.
\]
由椭圆方程
\[
y_0^2=16\left(1-\frac{x_0^2}{25}\right)
=16\left(1-\frac59\right)=\frac{64}{9},
\]
所以切点为
\[
\boxed{P\left(\pm\frac{5\sqrt5}{3},\ \pm\frac83\right)}
\]
（两个正负号可独立选择，共四个对称点）。
\textbf{进一步：} $r_1+r_2=10,\ r_1r_2=20$，故两焦半径为 $5+\sqrt5,5-\sqrt5$（不指定先后顺序）。

\medskip
\textbf{\textcolor{CExample}{例 3｜条件陷阱：原文件的数值为什么无解？}}

\textbf{题目：}椭圆 $\dfrac{x^2}{25}+\dfrac{y^2}{16}=1$，有切线距中心 $d=2$。求 $\sqrt{r_1r_2}$。

\textbf{辨析：}椭圆任何切线都满足 $b\le d\le a$，即 $4\le d\le5$。
$2<4$，故\textbf{不存在这样的切线}，原题\textbf{无解}。
若只机械套公式会得 $\sqrt{r_1r_2}=20/2=10$，
但由 $r_1+r_2=10$ 与均值不等式，
\[
\sqrt{r_1r_2}\le\frac{r_1+r_2}{2}=5,
\]
形成矛盾。这是非常好的\textbf{考场反检验}。

\medskip
\textbf{\textcolor{CExample}{例 4｜几何角度：不写坐标也能算}}

\textbf{题目：}椭圆半长轴 $a=5$、半短轴 $b=3$，点 $P$ 的焦点夹角满足 $\cos(\angle F_1PF_2/2)=\dfrac45$。求 $d$ 和 $r_1r_2$。

\textbf{解：}几何推论直接给
\[
d=a\cos\frac\theta2=5\cdot\frac45=\boxed{4},
\]
\[
r_1r_2=\frac{b^2}{\cos^2(\theta/2)}
=\frac9{16/25}=\boxed{\frac{225}{16}}.
\]
检验：$r_1r_2d^2=(225/16)\cdot16=225=a^2b^2$。
\textbf{启发：}题目出现焦点三角形的角时，用几何法比设 $P(x_0,y_0)$ 快。

\medskip
\textbf{\textcolor{CExample}{例 5｜迁移：焦点距切线之积（知识图谱连接）}}

\textbf{题目：}同一椭圆 $a=5,b=3$，点 $P$ 的切线到中心距离 $d=4$，分别记 $h_1,h_2$ 为两焦点到该切线的距离。求 $h_1h_2$ 和 $r_1r_2$。

\textbf{解：}
\[
\boxed{h_1h_2=b^2=9},\qquad
\boxed{r_1r_2=\frac{a^2b^2}{d^2}
=\frac{225}{16}}.
\]
校验： $h_i/r_i=d/a=4/5$，所以 $h_1h_2=(16/25)r_1r_2=9$。
\textbf{迁移：}原点距切线与焦点距切线是\textbf{不同的距离}，但可以通过 $O$ 为焦点中点及光学性质联系起来。
\end{examplebox}
